\section{Experiments and Results}
In this section, we conduct extensive experiments to evaluate the performance of the proposed Cross-Modality Semantic Alignment Transformer (CMSA-Trans) for the zero-shot fault diagnosis of rotating machinery. The evaluation focuses on the ability of the model to generalize to unseen fault categories by using Large Language Model (LLM)-generated semantic descriptions.

\subsection{Experimental Setup}
To validate the effectiveness of the proposed method, we use two widely recognized mechanical fault diagnosis benchmarks: the Case Western Reserve University (CWRU) bearing dataset and the South East University (SEU) gearbox dataset. These datasets provide a full evaluation across different rotating components (bearings vs. gears) and varying levels of background noise and mechanical complexity.

The CWRU bearing dataset is a gold-standard benchmark in the domain of rotating machinery health monitoring, consisting of vibration measurements acquired under a wide spectrum of drive-end and fan-end motor-driven states. In this study, we use signals collected from a 2-hp Reliance Electric motor-driven mechanical system across four distinct operational load conditions (0, 1, 2, and 3 hp) and varying rotating frequencies. The primary data segments are sampled at a rate of 12 kHz, which captures the fine-grained high-frequency transient impacts characteristic of industrial faults. Our experimental corpus includes 3,500 signal segments, where each segment is meticulously labeled according to its fault diameter (0.007, 0.014, 0.021, or 0.028 inches) and localized component position.

For our specialized zero-shot experimental configuration, we partition the fault categories into seen and unseen sets following a rigorous protocol to simulate real-world failure detection scenarios. Specifically, health states in the seen training set include the normal baseline (N), ball-induced defect (B, 0.007'' diameter), and outer race fault (OR, 0.014'' diameter) under a steady-state condition of 1797 rpm. In contrast, inner race faults (IR, 0.021'' diameter) and complex multi-component fault combinations (such as a 0.028'' ball failure) are defined as unseen categories, which are deliberately withheld during the training phase to test the zero-shot inference capability of the model. The data splitting strategy follows an 80/20 ratio for seen classes, where 80\% (2,000 samples) are used for optimization and 20\% (500 samples) are used for validation. The 1,000 samples belonging to unseen classes are reserved exclusively for the final zero-shot evaluation phase. This disjoint classification scheme mandates that the model transfers the learned physical knowledge from observed classes through the high-level semantic bridge of LLM descriptions.

The SEU gearbox dataset, sourced from the South East University dynamics laboratory, provides a more challenging testbed by introducing complex mechanical interactions within a simulated industrial gearbox architecture. This system features a multi-stage transmission involving a motor, a planetary gear reducer, a parallel shaft gearbox, and a magnetic powder brake for torque loading. The presence of overlapping vibration signals from diverse gear pairs creates a significant cross-component noise environment. Data segments are sampled at 5120 Hz across rotating speeds of 20 Hz and 30 Hz under varying operational loads of 0 V and 7.3 V. The dataset contains 4,000 segments, each representing distinct health conditions under different load profiles. Within this dataset, we characterize five prevalent gearbox health states: the normal monitoring state, the chipped tooth gear failure, the worn tooth condition, the catastrophic broken tooth state, and the root-based tooth fault. For our cross-modality zero-shot alignment experiments, we designate the Normal, Broken, and Chipped states as the seen classes (consisting of 2,400 training/validation samples), while the Worn and Root faults are reserved as the unseen target categories (1,600 test samples). Technical text-based semantic descriptors for these specific faults are generated by querying the GPT-4 LLM with technical prompt engineering. These prompts use expert-level diagnostic concepts—such as "non-stationary modulation patterns," "circular pitch error impacts," and "energy leakage in the mesh frequency sidebands"—to provide the system with a detailed technical map of each zero-shot target class. By using such high-fidelity linguistic information, the CMSA-Trans can better distinguish between geometrically similar failures, such as worn and chipped teeth, in the absence of labeled signals.

Signal preprocessing is a high-priority phase for ensuring the reliability and efficiency of the transformer-based architectural encoder. The raw vibration signals, often contaminated by industrial environment interference, are first segmented using an overlapping sliding window approach. We employ a fixed window size of 1024 data points combined with an overlap of 512 samples. This segmentation strategy is important. It ensures sufficient data diversity for model generalization. It also preserves critical transient features—such as periodic impacts and impulsive components—that are necessary for accurate zero-shot fault identification. This overlap ensures that transient events occurring at the boundaries of windows are captured in at least one segment, preventing information loss. Subsequently, each segment undergoes Gaussian Z-score normalization, which is key to eliminate the bias from stationary DC components and minimize the influence of varying operating conditions, including fluctuations in motor load and rotating speed. Mathematically, for a raw signal $x[n]$, the normalized signal $\hat{x}[n]$ is computed as $(x[n] - \mu) / \sigma$, where $\mu$ and $\sigma$ represent the mean and standard deviation of the current window. To further enhance the spectral representation, we apply a Hanning window to reduce spectral leakage before feeding the signal into the multi-scale 1D-CNN feature extractor. The resulting normalized temporal sequences, mathematically represented as $X \in \mathbb{R}^{1024 \times 1$, are used as the direct input tensors. This setup yields a total of 2500 samples for the seen classes and approximately 1000 samples for the unseen categories in each dataset. This process provides a balanced yet rigorous environment for zero-shot diagnostic performance assessment, challenging the model to bridge the gap between distinct modalities under varying operational constraints. In terms of data partitioning and distribution, 80\% of the seen class samples are allocated for model training, while the remaining 20\% serve as the validation set to monitor the convergence of the cross-modal semantic alignment loss and prevent overfitting to the source domain.

\subsection{Implementation Details}
The proposed CMSA-Trans framework is implemented using the PyTorch 2.1 deep learning library, using its dynamic graph capabilities for multi-modal fusion. The architecture is modularly designed and optimized for GPU parallelization, consisting of a custom three-layer 1D-CNN backbone for deep local feature extraction, followed by a four-layer Transformer encoder block integrated with eight multi-head self-attention units. Each individual CNN layer uses a specialized kernel size of 3 and a stride of 1. This configuration effectively reduces the temporal dimensionality of the vibration waveform while simultaneously expanding the feature channel space to capture higher-level abstractions. The final embedding dimension $d_{model}$ within the latent space of the Transformer is set to 256. This specific dimensionality was chosen after extensive preliminary testing, as it provides an optimal balance between representational expressiveness and computational efficiency.

For the language processing branch, we use the pre-trained BERT-large model with Whole Word Masking (WWM). The LLM-generated technical fault semantics are processed by the BERT encoder to produce 1024-dimensional semantic embedding vectors. These linguistic representations are then mapped into the shared 256-dimensional joint embedding space through a dedicated two-layer feed-forward neural network (FFN), which acts as a semantic projection head. The model is optimized using the Adam optimizer with a weight decay coefficient of $1e-4$. The initial learning rate is set to $0.001$, and we employ a Cosine Annealing scheduler to dynamically adjust the learning rate, allowing for stable convergence over a 200-epoch training trajectory. We use a batch size of 64 to maintain high gradient fidelity during the stochastic optimization process. All training and inference tasks are performed on a workstation equipped with an NVIDIA GeForce RTX 3090 GPU (24GB VRAM) and an Intel Core i9-12900K processor. The total training duration for a single complete dataset is approximately 55 minutes. The important hyperparameter $\lambda$ in the utility function Eq. 9, which regulates the sensitivity and trade-off between the cross-modality alignment loss and the source domain classification accuracy, is empirically optimized and fixed at 0.5. To eliminate the stochastic variability of random weight initialization and ensure the statistical reliability of the findings, every experiment is executed 10 times independently. The mean accuracy values and their corresponding standard deviations are then documented as the performance metrics for this study.

\subsection{Zero-Shot Recognition Performance}
The zero-shot diagnostic performance of the CMSA-Trans is rigorously benchmarked against five leading baseline methods to highlight its superiority in cross-modality alignment. These baselines include: (1) WDCNN \cite{zhang2017new}, a classic CNN-based model adapted for ZSL via attribute matching; (2) ZSL-GAN \cite{xian2019fvaegan}, a generative adversarial network designed to synthesize unseen class features to bridge the domain gap between observed and new faults; (3) Deep Multi-modal ZSL (DMZSL) \cite{schonfeld2019generalized}, which uses a shared subspace for vision and language to enable knowledge transfer; (4) Attribute-based ZSL (AZSL) \cite{lampert2014attribute}, relying on manual engineer-defined fault attributes; and (5) Bi-LSTM-ZSL \cite{chen2019bearing}, which captures temporal dependencies but lacks specialized semantic alignment layers. As summarized in Table~\ref{tab:comparison}, CMSA-Trans demonstrates a marked performance advantage across all fault categories \cite{li2022zero, lu2017deep}.

On the CWRU bearing dataset, CMSA-Trans achieves an average zero-shot accuracy of 89.4\% for the unseen fault categories. This represents a 4.2\% improvement over the runner-up method (ZSL-GAN). While GAN-based methods often struggle with instability and generate noisy feature distributions that lack physical consistency \cite{lei2020applications}, the Transformer-based CMSA-Trans architecture provides a stable and structured mapping between the vibration signal patches and the LLM semantic tokens. For the "Inner Race Fault" (unseen), CMSA-Trans achieves 91.2\% accuracy, greatly outperforming AZSL (76.8\%) and WDCNN (71.3\%). This improvement is due to the ability of the LLM to describe the "higher-order harmonics" and "envelope modulation" typically seen in inner race defects, which are often overlooked in manual attribute sets. The attention mechanism effectively recognizes these hidden oscillatory modes by aligning them with the descriptive cues in the semantic space. Deep architectures like WDCNN often fail to capture these subtle signatures. Their fixed kernels are insensitive to the long-range temporal periodicity of inner race impacts \cite{hochreiter1997long}. In contrast, our model treats each impact as a distinct semantic token, allowing for a more granular identification process.

On the more complex SEU gearbox dataset, the advantages of the semantic alignment strategy are more pronounced. CMSA-Trans achieves a zero-shot accuracy of 82.5\% on this dataset. Traditional attribute-based methods (AZSL) drop to 65.4\%. The performance gap shows that the rich, descriptive technical semantics provided by LLMs act as a superior bridge for knowledge transfer compared to sparse binary attributes. Specifically, for the "Broken Tooth" category, CMSA-Trans maintains 88.7\% accuracy, which is 15.2\% higher than Bi-LSTM-ZSL. Root faults exhibit unique spectral signatures in the sideband frequencies. These are explicitly detailed in our LLM prompts but are ignored by LSTM-based sequence models that only look at global temporal patterns \cite{zhang2021attention}. Fig.~\ref{fig:accuracy_bar} presents the confusion matrix for the unseen categories. It shows that the model rarely confuses the broken tooth fault with the normal state, thanks to the distinct semantic prototypes provided by the language model. The high precision and recall across unseen categories demonstrate that CMSA-Trans is highly effective at generalizing to new mechanical health states without requiring any labeled signal data during training, showcasing its potential for practical prognostic and health management (PHM) applications \cite{jia2016deep}.

\subsection{Ablation Study}
To further dissect the contribution of each component within the CMSA-Trans framework, we conduct a systematic ablation study on both the CWRU and SEU datasets. We evaluate four distinct model configurations: (1) CMSA-Trans-Full: the proposed architecture with all components active; (2) CMSA-Trans-NoLLM: uses traditional 20-dimensional manual attribute vectors instead of 1024-dimensional LLM-generated semantic embeddings; (3) CMSA-Trans-NoAttn: replaces the cross-modality attention alignment module with global average pooling followed by a concatenation fusion of signal and text features; (4) CMSA-Trans-Base: replaces the Transformer encoder with a standard 1D-ResNet structure, thereby removing the self-attention mechanism for temporal dependency modeling.

The quantitative results across both datasets highlight a consistent trend. For the CWRU dataset, CMSA-Trans-Full achieves 89.4\% accuracy, while CMSA-Trans-NoLLM, CMSA-Trans-NoAttn, and CMSA-Trans-Base achieve 81.6\%, 83.9\%, and 85.1\%, respectively. On the more challenging SEU dataset, CMSA-Trans-Full maintains a high 82.5\%, whereas the variants drop markedly to 74.7\%, 77.0\%, and 78.4\%. As shown in Table~\ref{tab:ablation}, the removal of LLM-generated semantics (CMSA-Trans-NoLLM) results in the most large performance degradation, with a 7.8\% drop in zero-shot accuracy on the SEU dataset. This confirms our hypothesis that the "semantic bottleneck" in traditional ZSL for fault diagnosis is caused by the low information density and subjective nature of manual attributes. LLMs provide a high-dimensional, continuous semantic space that better preserves the transitive relations between different fault types by encoding linguistic relationships.

Also, the exclusion of the cross-attention mechanism (CMSA-Trans-NoAttn) leads to a 5.5\% accuracy loss. This indicates that local alignment between specific signal segments (e.g., transient impacts) and semantic keywords is key for fine-grained fault recognition. Simple global fusion fails to capture the intricate temporal-semantic correlations. Without the attention layer, the model attempts to map the entire averaged signal vector to the entire semantic prototype, losing the ability to associate specific pulses with specific textual descriptors. Finally, the comparison with CMSA-Trans-Base shows that the self-attention layers of the Transformer are better at capturing the long-range periodic nature of rotating machinery vibrations compared to the local receptive fields of standard residual convolutions. The ResNet-based backbone (CMSA-Trans-Base) struggles with signals collected under varying load conditions, where the fault impacts vary in frequency and phase. The ability of the Transformer to attend to all time steps simultaneously allows it to remain invariant to these phase shifts, which is essential for consistent zero-shot alignment across different operating speeds.

\subsection{Parameter Sensitivity and Robustness}
We conduct a sensitivity analysis to determine the optimal configuration of the latent embedding space and evaluate the robustness of the model against measurement noise. Fig.~\ref{fig:sensitivity} illustrates the variation in zero-shot accuracy as the hidden embedding dimension $d_{model}$ is tuned from 64 to 512. We observe that increasing the dimension from 64 to 256 yields a sharp improvement in performance as the model gains the capacity to represent complex cross-modal relationships between the vibration and semantic domains. However, a further increase to 512 results in a slight decrease in validation accuracy, likely due to optimization challenges in the higher-dimensional space. So, 256 is selected as the optimal dimension for our final architecture.

The robustness of the proposed CMSA-Trans is an important factor for its deployment in harsh industrial environments. To test this, we inject Additive White Gaussian Noise (AWGN) into the raw test signals to simulate different Signal-to-Noise Ratios (SNR), ranging from -10 dB to 10 dB. As shown in the robustness curve in Fig.~\ref{fig:snr_curve}, CMSA-Trans maintains a relatively stable performance even in low SNR conditions. At SNR = 0 dB, CMSA-Trans achieves 75.6\% accuracy, which is 12.3\% higher than the performance of WDCNN under the same conditions. This superior noise resistance is attributed to the multi-scale CNN backbone that extracts noise-invariant features and the Transformer-based attention mechanism that acts as a temporal filter, focusing on the quasi-periodic fault impacts. These results emphasize that CMSA-Trans is not only accurate in laboratory conditions but also reliable for actual field diagnosis where noise interference is inevitable.

\subsection{Visual Analysis}
In this subsection, we provide a visual interpretation of the alignment mechanism to explain why the model generalizes effectively to unseen classes. We first use the t-Distributed Stochastic Neighbor Embedding (t-SNE) algorithm to project the 256-dimensional features from the alignment layer into a 2D space. As illustrated in Fig.~\ref{fig:tsne}, the signal features of the unseen classes (Inner Race and Combined Faults) form distinct, high-density clusters that are spatially adjacent to their corresponding LLM-generated semantic prototypes. The spatial distance between the clusters in the embedding space mirrors the semantic distance between their textual descriptions, validating that the model successfully learns a physically meaningful semantic manifold.

Also, we visualize the attention scores between the input vibration signal patches and the semantic tokens to assess the explainability of our zero-shot predictions. As shown in the attention heatmaps (Fig.~\ref{fig:attn_map}), when presented with an unseen inner race fault signal, the attention heads consistently assign high weights to semantic tokens such as "periodic peak," "carrier frequency," and "bearing housing resonance." This indicates that the CMSA-Trans is able to recognize the physical characteristics within the raw waveform and match them with the conceptual keywords provided by the LLM. This evidence suggests that the cross-modality semantic alignment transformer establishes a deep logical connection between the physical laws of machinery failure and natural language descriptions. Such a mechanism provides a transparent and strong foundation for autonomous fault diagnosis in smart manufacturing systems. This allows engineers to verify the diagnostic logic of the model via human-understandable text alignment.
